Name based audits often treat evidence of a social association as evidence of downstream discrimination. We test whether that inference is warranted using Chilean surnames as controlled socioeconomic probes. Phase II evaluates eight frozen model/provider cells on 1,032 prompts each, for 8,256 verified responses, separating forced latent association from matched consequential decisions across academic selection, hiring, research fellowship selection, and legal aid intake. Elite coded surnames received higher forced high status probability mass than common surnames in seven of eight models and higher mass than rare frequency controls in all eight models. Yet elite minus common decision effects were close to zero for most systems: five models were statistically equivalent within a predeclared $\pm0.10$ standard deviation margin, while the remaining three were either imprecise or borderline rather than showing a consistent elite advantage. Association strength did not predict decision leakage across models ($r=0.201$, $p=0.633$) or across frozen surname pair by model cells ($r=0.065$, $p=0.565$). The results show that latent social association and consequential treatment are empirically distinct constructs and should be evaluated separately rather than inferred from one another.